\section{Introduction}
\label{sec:introduction}

The earlier chapters of this thesis treated the Gromov--Wasserstein
(GW) distance as a static object: a geometry-aware divergence
between metric--measure spaces \citep{memoli_2022_distance},
defined purely in terms of the intra-domain pairwise-distance
matrices of the two sides, solved either exactly, through its
entropic relaxation
\citep{peyre_2016_gromov, solomon_2016_entropic} or fused-features variant \cite{vayer_2019_fused} and deployed \emph{post hoc} to compare or align fixed representations. This
chapter changes the role of GW. Rather than measuring how well
two given embedding sets can be put into correspondence, we ask
whether the same divergence is useful as a \emph{differentiable
regularizer} that shapes the representations of a deep generative
model \emph{during training}. The informal intuition is that 
if two real-world phenomena are similar in one modality, the
corresponding signals should be similar in the other. Pairs of
clips that share a sound source ought to lie close in audio space
to the same extent that they lie close in video space, and
isolated visual events should produce isolated acoustic events at
related temporal scales. This is a statement about
\emph{relational geometry} rather than about pointwise
correspondence, and the Gromov--Wasserstein distance is a natural way of formalising it.

The reason GW has accumulated traction across disparate fields is
precisely that it operates on intra-domain pairwise relations
rather than on features, and is therefore invariant to any
isometric transformations of either side. Two domains can live
in spaces of different dimensions, with different units, different
encoders, and even different topologies; GW only requires that
each side comes with its own notion of pairwise distance. In shape
analysis and computer graphics, this invariance is what allows
entropic GW to find stable soft correspondences between surfaces
and point clouds embedded in incomparable spaces, without
landmarks or a common parametrisation
\citep{peyre_2016_gromov, solomon_2016_entropic}. In vision and
language, the GOT framework of \citet{chen_2020_graph} adds a
fused Wasserstein--GW penalty to image--text retrieval, visual
question answering, captioning, machine translation, and text
summarisation, with the GW component contributing a structural
lift on top of pointwise attention; \citet{lu_2025_crossmodal}
report the same pattern when transferring linguistic knowledge
into a CTC speech recogniser, and \citet{chandra_2025_realign}
obtain double-digit improvements on procedural-video alignment
using a partial fused GW loss. In multimodal representation
learning, \citet{gong_2022_gromov} go further and show that
\emph{pure} GW is able to align modalities even when no paired
data are available, by matching their pairwise-distance structure
alone. 

In the comparative analysis of neural representations,
\citet{sasaki_2025_unsupervised} use GW to recover correspondences
between brains, species, and artificial networks that supervised
representational similarity analysis flattens, precisely because
GW does not require a known stimulus-to-stimulus map. And in
structured-data learning, the fused Gromov--Wasserstein extension
\citep{vayer_2019_optimal, vayer_2019_fused} compares attributed
graphs jointly through node features (Wasserstein) and edge
geometry (GW), outperforming graph kernels and graph convolutional
networks on standard benchmarks.

\cite{zhou2025fusedgromovwassersteincontrastivelearning} have implemented a Fused-Gromov-Wasserstein (FGW) regularization term along with other alignment losses, showing robust generalization on enzyme-reaction benchmarks and enzyme screenings compared to the state-of-the-art. 

The underlying lesson of these works is that the internal geometric structure of the different collections of objects can be leveraged as a complementary signal to point-wise associations \cite{chen_2020_graph, lu_2025_crossmodal, chandra_2025_realign, gong_2022_gromov, sasaki_2025_unsupervised, vayer_2019_optimal}.

The goal of this section is to assess whether the addition of Gromov-Wasserstein regularizations can improve on the downstream generation tasks. For consistency with the previous experiments in this thesis We test it on video-to-audio cross-modal generation. 
We argue that on a distributional level, similar visual sample representations will also share similar audio embeddings. More informally, the idea can be states as "similar looking things sound similarly". 
A barking dog and a meowing cat occupy nearby regions of visual semantic space, and their characteristic sounds likewise occupy nearby regions of audio semantic space; clips of
splashing water cluster together visually and acoustically; impact
events at a given temporal scale produce sounds at related
spectral scales. A video-to-audio (V2A) model that respects this
relational structure should, all else equal, generate audio whose
intra-batch geometry mirrors the intra-batch geometry of the
conditioning video, property that no current training objective enforces explicitly. 

Modern flow-matching V2A systems
\citep{wang_2024_frieren, cheng_2025_mmaudio}, like their
text-to-audio cousins \citep{guan_2024_lafma}, build on the
simulation-free flow-matching objective of
\citet{lipman_2023_flow, albergo_2023_building, liu_2023_rectified}
and inject video information through cross-attention, adaLN-style modulation and frame-aligned
conditioning. The loss itself is a per-token regression that does
not compare the geometry of the two streams. In the language of
optimal transport, the training objective is shaped by
Wasserstein-style (feature-level) signals, while the
Gromov--Wasserstein-style (relational) signal is left implicit, to
be discovered by the network if and when it helps.

To the best of our knowledge, no published work integrates a
Gromov-Wasserstein objective into a generative cross-modal
audio-video model. Existing cross-modal GW results target
discriminative tasks retrieval, captioning, ASR, segmentation,
clustering --- or self-supervised pretraining; the closest
generative analogues \citep{wang_2024_inverse, rho_2025_lavcap, you_2025_mover} use
only the Wasserstein component and do not exploit the
relational-geometry signal that distinguishes GW from plain OT.

In the methodology section (\S\ref{sec:methodology}) we test a flow-matching V2A backbone augmented
with a per-mini-batch entropic-GW penalty between intra-video and
intra-audio pairwise distance matrices. We run experiments with both a fixed and scheduled weighted $\lambda_{\mathrm{GW}}(s)$ coefficient. We deliberately keep the backbone fixed and treat it purely as the basis to test our regularizer ablations, so that any improvement or degradation is attributable to the relational signal rather than to architectural change. 

The chapter contributes by integrating a Gromov-Wasserstein regularizer, in a fashion similar to the previously discussed works (e.g. \citet{chen_2020_graph, zhou2025fusedgromovwassersteincontrastivelearning} into a flow-matching cross-modal
generative objective as novel contribution.  
 

Second, an ablation across five
representation variants that isolates \emph{where} in the network the
relational signal should be injected. 

Third, a schedule for the
trade-off coefficient $\lambda_{\mathrm{GW}}(s)$ that resolves an
engineering question left open by prior fused-OT work
\citep{chen_2020_graph, vayer_2019_optimal}, where the mixing
parameter is typically held fixed. 

Fourth, a spectral
$-\log\det$ Gram penalty \citep{ozsoy_2022_corinfomax} adapted
from non-contrastive self-supervised learning
\citep{bardes_2022_vicreg, zbontar_2021_barlow,
ermolov_2021_whitening} to forbid the rank-1 collapse mode that
pure GW alignment between cross-batch representations would produce when learning projections.
Fifth, an empirical evaluation on VGGSound asking whether the
relational signal contributes any improvement on top of the
already strong pointwise conditioning of a modern V2A backbone.


The use of entropic GW as a per-batch loss is justified by the
recently established statistical theory:
\citet{zhang_2024_gromov} show that entropic GW achieves the
parametric $n^{-1/2}$ sample-complexity rate independent of the
ambient dimension, and \citet{rioux_2023_entropic} prove
convergence of the entropic-GW solver together with a
central-limit theorem for the empirical estimator. The remaining
design tensions --- the cubic per-iteration cost of standard
entropic GW versus the linear-time low-rank variant of
\citet{scetbon_2022_lineartime} or the sliced variant of
\citet{vayer_2019_sliced}, and the equal-mass assumption versus
the unbalanced formulation of \citet{sejourne_2021_unbalanced}
--- are taken up in the methodology and discussion sections.

The remainder of this chapter is organised as follows. Section
\ref{sec:related_work} situates the contribution among prior
cross-modal OT/GW work and modern V2A systems. Section
\ref{sec:methodology} introduces the conditional flow-matching
backbone (concretely realised by the open-source MMAudio system
\citep{cheng_2025_mmaudio}, on top of which our regularizer is
implemented), defines the five GW representation variants, the
fused Gromov--Wasserstein objective, the spectral anti-collapse
penalty, and the $\lambda_{\mathrm{GW}}$ schedule. Sections
\ref{sec:results} and \ref{sec:discussion} report and interpret
the resulting evidence.
